% Slide — Os dois defeitos, em detalhe
\begin{frame}{Os dois defeitos que motivaram a reescrita}
  \tabstyle\footnotesize
  \begin{tabular}{@{}p{0.16\linewidth}p{0.40\linewidth}p{0.40\linewidth}@{}}
    \toprule
    & \antes & \depois \\
    \midrule
    \textbf{PPO} &
      Ator atualizado somando a vantagem $A_t$ diretamente aos logits-alvo e
      treinado por \textbf{erro quadrático}. Isso não é o gradiente da
      \textit{surrogate function} recortada do PPO. &
      Gradiente correto por autograd sobre a razão de verossimilhanças
      $r_t=\pi_{\text{novo}}/\pi_{\text{antigo}}$, com \textbf{clipping}
      explícito e \textbf{GAE}$(\lambda)$. \\
    \addlinespace
    \textbf{Aptidão GA/SA/PSO/DE} &
      Soma ponderada fixa $\lambda_1\,\mathrm{NS}-\lambda_2\,\mathrm{CTI}$.
      O ponto de operação depende da \textbf{escala do CTI}, que varia até
      124$\times$ entre a maior e a menor loja do recorte piloto. &
      Formulação \textbf{restrita} (Eq.~4.2 da proposta):
      $\min \mathrm{CTI}$ s.a. $\mathrm{NS}\geq\alpha_{\min}$, com
      penalidade \textbf{relativa ao próprio CTI} -- adimensional, comparável
      entre séries. \\
    \bottomrule
  \end{tabular}
  \vspace{0.8em}

  \begin{alertblock}{Por que isso importa para os resultados já reportados}
    \small A degeneração observada no Experimento 1 (frequência de pedidos
    FP $=0{,}98$ para a política PPO -- pedir quase todo ciclo) é
    \highlight{consistente com o defeito do gradiente}, não com uma
    limitação real do método PPO.
  \end{alertblock}

  \note{
    O defeito do PPO é sutil de detectar olhando só para o código, mas
    óbvio olhando para o resultado: um agente PPO bem treinado não deveria
    pedir em 98\% dos ciclos -- isso é comportamento de política quase
    aleatória, não de política otimizada. A causa raiz é que a atualização
    do ator não seguia o gradiente do objetivo clipped surrogate do PPO; ela
    empurrava os logits na direção da vantagem e depois ajustava por MSE,
    o que é uma heurística sem garantia de melhoria monotônica. O segundo
    defeito, da aptidão das meta-heurísticas, é mais estrutural: como o
    custo varia por ordens de grandeza entre séries de volumes diferentes,
    o mesmo par de pesos (lambda1, lambda2) que funciona bem numa loja de
    alto giro pode privilegiar demais o serviço numa loja pequena, ou
    demais o custo numa loja grande -- contaminando qualquer comparação
    entre políticas feita com a mesma configuração de pesos.
  }
\end{frame}
